\chapter*{MỞ ĐẦU}

Kiểm tra chất lượng là khâu bắt buộc trong sản xuất công nghiệp: sản phẩm có vết xước,
nứt, nhiễm bẩn hay thiếu chi tiết phải bị loại trước khi đến tay người dùng. Kiểm tra bằng mắt người chậm, phụ thuộc
kinh nghiệm và dễ bỏ sót khi sản lượng lớn hoặc lỗi nhỏ. Thị giác máy tính có thể tự động hóa khâu này, nhưng gặp một
khó khăn về dữ liệu: ảnh sản phẩm lỗi hiếm, đa dạng và khó lường trước, còn việc khoanh vùng lỗi trên từng ảnh để làm
nhãn tốn nhiều công sức, trong khi mô hình phân loại có giám sát chỉ nhận ra những kiểu lỗi đã thấy. Ngược lại, ảnh sản phẩm tốt có sẵn với số lượng lớn. Vì vậy, hướng học không giám sát, tức chỉ học từ ảnh tốt rồi
coi mọi sai khác đáng kể là bất thường, phù hợp hơn với điều kiện thực tế~\cite{bergmannMVTecADComprehensive2019}.

Đề tài xây dựng và đánh giá ba phương pháp phát hiện và định vị bất thường không giám sát trên
ảnh sản phẩm công nghiệp, nhằm trả lời ba câu hỏi: phương pháp nào phát hiện và khoanh vùng lỗi chính xác nhất;
phương pháp nào phù hợp với loại sản phẩm nào (bề mặt hay vật thể, vị trí cố định hay xoay); và mỗi phương pháp đánh
đổi giữa độ chính xác với thời gian, bộ nhớ ra sao.

Ba phương pháp được chọn đại diện cho hai hướng tiếp cận. Autoencoder, do nhóm tự xây
dựng, học tái tạo ảnh từ đầu và dùng sai số tái tạo làm tín hiệu bất thường.
PaDiM~\cite{defardPaDiMPatchDistribution2021} và PatchCore~\cite{rothTotalRecallIndustrial2022} dùng đặc trưng của
mạng tích chập đã học trước trên ImageNet, lần lượt mô hình hóa đặc trưng bình thường bằng phân phối Gauss theo vị
trí và bằng một bộ nhớ các vector đặc trưng. Dữ liệu là đủ 15 loại sản phẩm của bộ MVTec
AD~\cite{bergmannMVTecADComprehensive2019}. Đề tài dừng ở mức đánh giá trên dữ liệu ảnh có sẵn, chưa triển khai trên
dây chuyền và chưa chọn ngưỡng quyết định cho sử dụng thực tế.

Về cách làm, nhóm tìm hiểu cơ chế của từng phương pháp qua bài báo gốc, cài đặt chúng theo cùng
một giao thức, chạy trên cả 15 loại và đánh giá bằng
image-AUROC cho nhiệm vụ phát hiện ảnh lỗi, pixel-AUROC cho nhiệm vụ định vị vùng lỗi.
Kết quả cho thấy PatchCore đạt image-AUROC trung bình 0,983 và pixel-AUROC 0,980, sát số liệu công bố; PaDiM
đạt 0,890 và 0,960 với thời gian xây dựng mô hình ngắn nhất; Autoencoder đạt 0,756 và 0,843. Như vậy, với
vài trăm ảnh tốt cho mỗi loại, dùng lại đặc trưng học trước hiệu quả hơn nhiều so với học tái tạo ảnh từ đầu.

Chương 1 trình bày cơ sở lý thuyết: học sâu, mạng tích chập, đặc trưng học trước, bài toán
phát hiện bất thường và cách đánh giá. Chương 2 trình bày bộ dữ liệu, giao thức đánh giá và cách triển khai ba
phương pháp. Chương 3 trình bày kết quả thực nghiệm và đánh giá.
